\subsection*{Хеш-функция.}

\Def Пусть $U$ - множество рассматриваемых объектов, тогда $h: U \rightarrow \{0, 1, \dots, m - 1\}$ называется хеш-функцией.

\subsection*{Коллизия.}

\Def Элементы $x, y \in U$, где $x \neq y$ образуют \textit{коллизию}, если $h(x) = h(y)$.

\subsection*{Хеш-таблица.}

Заведем массив размера $|h(U)|$ и будем класть булев флаг в ячейку $h(x)$, если элемент x 
есть в массиве. 

Минусы:
\begin{itemize}
    \item Огромные затраты памяти $\rightarrow$ можно завести массив размера $m$ и класть элемент по индексу $h(x) \% m$.
    \item Коллизии.
\end{itemize}

\subsection*{Операции}
\begin{itemize}
    \item Insert - добавляет элемент.
    \item Delete - удаляет элемент.
    \item Contains (Random access) - проверка на наличие элемента в таблице.
\end{itemize}

\subsection*{Хеш-таблица на цепочках.}

\Def Бакет (bucket) - нечто, хранящее все элементы, образующие коллизию.

Вместо обычной хеш-таблицы заведем хеш-таблицу на цепочках. Вместо булевых флагов в ячейках $h(x)$ будем хранить цепочки (бакеты) с элементами которые образовали коллизию. То есть получим массив цепочек, где все операции работают за длину цепочки.

\subsection*{Время работы.}

Если соблюдать ограничение на \textit{load factor}, то сложность всех операций составит $\mathcal{O}(1)$ в среднем. 
\newline 
$\rightarrow$ Вставка, удаление и поиск в среднем за $\mathcal{O}(1)$.

\Note \textit{load factor} - коэффициент загруженности (подробнее в след. билете).

\pagebreak